\documentclass{mcom-l}
\issueinfo{00}{0}{Xxxx}{XXXX}

\usepackage{amsmath,amssymb,mathtools}
\usepackage{booktabs}
\usepackage{graphicx}
\usepackage{url}
\usepackage[hidelinks]{hyperref}

\newtheorem{theorem}{Theorem}[section]
\newtheorem{lemma}[theorem]{Lemma}
\newtheorem{corollary}[theorem]{Corollary}
\theoremstyle{definition}
\newtheorem{definition}[theorem]{Definition}
\newtheorem{example}[theorem]{Example}
\theoremstyle{remark}
\newtheorem{remark}[theorem]{Remark}

\newcommand{\Ftwo}{\mathbb{F}_2}
\newcommand{\poly}{\Ftwo[x]}
\newcommand{\Dmin}{\Delta_{\min}}
\newcommand{\Wfb}{W_{\mathrm{fb}}}
\newcommand{\drafttodo}[1]{%
  \par\smallskip\noindent\textbf{[TODO]} \emph{#1}\par\smallskip}
\newcommand{\prooftodo}{%
  \par\smallskip\noindent\textbf{[Proof TODO.] }%
  Complete the formal proof and audit all boundary cases.\par\smallskip}
\newcommand{\placeholderbox}[1]{%
  \begin{center}
    \fbox{\parbox{0.90\linewidth}{\centering\textbf{Placeholder.} #1}}
  \end{center}}

\title[Generalized Suwako Reduction]{Generalized Suwako:\
Sparse Modular Reduction with Logarithmic Feedback Depth}

% TODO: Replace all author metadata before circulation.
\author{Author Name}
\address{Institution and postal address}
\email{author@example.org}
\subjclass[2020]{Primary 11T06; Secondary 12Y05, 68W10, 68W40}
\keywords{binary polynomials, modular reduction, finite fields,
sparse polynomials, Frobenius map, parallel algorithms}

\begin{document}

\begin{abstract}
Polynomial modular reduction is a basic operation in binary finite-field
arithmetic and in computational-algebra workloads based on repeated Frobenius
powers.  For a sparse modulus, conventional top-down folding retains sparse
shift/XOR work, but its dependency chain can be long when a nonleading term is
close to the leading term.  We study a different schedule obtained by
expressing the high-part recurrence as the inverse of a nilpotent feedback
operator.  In characteristic two, this inverse factors into Frobenius powers
whose individual supports remain controlled by the original tap set.  The
result is a matrix-free reduction method with an exact feedback-stage count
and a scheduled-work expression governed by the complete tap geometry.

\drafttodo{Replace this paragraph with the final proved complexity statement,
matched reduction results, and one end-to-end repeated-squaring,
irreducibility-testing, or factorization result.  Include the platform,
parameter range, strongest baselines, and losing regions.}
\end{abstract}

\maketitle

\section{Introduction}
\label{sec:introduction}

Arithmetic in a polynomial-basis representation of a binary extension field
reduces products and squares modulo a monic polynomial.  The same operation
appears in polynomial factorization, irreducibility testing, and searches for
sparse irreducible or primitive polynomials.  In these settings the modulus is
public and often fixed for many reductions, but the useful degree, support,
and setup-amortization regimes vary substantially.

When the nonleading part is sparse, top-down folding uses little setup and can
be extremely effective.  Its work and latency, however, are controlled by
different geometric features.  Hamming weight controls how many shifts and
XORs one fold introduces, whereas the distance from the leading term to the
nearest internal tap controls how far feedback can propagate.  Equal-degree,
equal-weight moduli may therefore induce very different dependency chains.

Barrett and Montgomery reduction avoid this particular serial recurrence, but
introduce multiplication kernels, modulus-dependent constants, and different
setup costs.  Reciprocal methods and dense linear maps expose additional
parallelism, but do not automatically preserve sparse support.  We ask whether
the sparse feedback chain can be shortened without first materializing a dense
reciprocal or reduction matrix.

For an input of degree below twice the modulus degree, reduction gives a
finite forced recurrence.  We write it as $(I+U)Y=H$ for a truncated shift
operator $U$.  Characteristic two gives
\begin{equation}
  (I+U)^{-1}=\prod_{k=0}^{r-1}(I+U^{2^k}),
  \qquad
  r=\left\lceil\log_2\frac{m}{\Dmin}\right\rceil,
  \label{eq:intro-factorization}
\end{equation}
once $U^{2^r}=0$.  Each Frobenius power doubles the original tap distances.
Mixed feedback paths arise through composition across stages rather than by
expanding a dense reciprocal inside one stage.

The intended contributions are: (i) a nilpotent-feedback formulation for any
monic binary modulus; (ii) a sparse characteristic-two factorization;
(iii) exact work--feedback-depth geometry; (iv) a matched implementation
study; and (v) evaluation in repeated modular squaring and one complete
computational-algebra workload.

\drafttodo{Items (i)--(iii) remain candidate theorem claims until the proof
and cost-model audit is complete.  Items (iv)--(v) remain planned until the
experiments in Sections~\ref{sec:evaluation} and~\ref{sec:consequences} are
complete.  Delete unsupported contributions.}

Correctness does not require irreducibility; sparsity is a performance
condition rather than an applicability condition.  We do not identify
feedback depth with Boolean gate depth, claim an unconditional improvement to
complete field multiplication, or infer a side-channel result from a public
fixed schedule alone.

Sections~\ref{sec:preliminaries}--\ref{sec:paradigms} fix the problem and
comparison models.  Sections~\ref{sec:operators}--\ref{sec:complexity} develop
the method and analysis.  Section~\ref{sec:related} states the prior-art
boundary.  Sections~\ref{sec:evaluation} and~\ref{sec:consequences} contain the
implementation and application evidence, and Section~\ref{sec:conclusion}
concludes.

\section{Preliminaries and Cost Models}
\label{sec:preliminaries}

Let
\begin{equation}
  g(x)=x^m+q(x)=x^m+\bigoplus_{t\in T}x^t\in\poly,
  \qquad T\subseteq\{0,\ldots,m-1\},
  \label{eq:modulus}
\end{equation}
be monic, and set $h=1+|T|$.  For $t\in T$, define
$\Delta_t=m-t$.  If $T\ne\varnothing$, let
\begin{equation}
  \Dmin=\min\{\Delta_t:t\in T\}.
\end{equation}
For $T=\varnothing$, the feedback operator below is zero and we set $r=0$.

Identify degree-below-$m$ polynomials with vectors in $\Ftwo^m$.  Let $N$ be
the truncated right shift
\begin{equation}
  (N^d y)_i=\begin{cases}y_{i+d},&0\le i<m-d,\\0,&m-d\le i<m,
  \end{cases}
  \qquad N^m=0.
\end{equation}

For $A\in\poly$ with $\deg A<2m$, write
\begin{equation}
  A=L+x^mH,
  \qquad \deg L,\deg H<m.
  \label{eq:low-high}
\end{equation}
For degree-below-$m$ input $Y$, define linear maps $V$ and $U$ by
\begin{equation}
  qY=V(Y)+x^mU(Y),
  \qquad
  U=\sum_{t\in T}N^{\Delta_t}.
  \label{eq:UV}
\end{equation}
Monicity is sufficient for a unique remainder; irreducibility is required only
when the quotient is intended to be a field.

We distinguish scheduled coefficient work, machine-word work, feedback depth,
bounded-fan-in gate depth, setup, and temporary space.  Scheduled coefficient
work is neither an instruction count nor a minimal XOR-circuit count.  Machine
latency and throughput are measured under a fixed platform and timed boundary.

\drafttodo{Freeze the word-operation model, representation, top-word masking,
allocation boundary, and setup-amortization policy.}

\section{Existing Reduction Paradigms}
\label{sec:paradigms}

Serial folding repeatedly eliminates high terms through the sparse support.
It has low setup, but a tap at distance $\Dmin$ can create on the order of
$m/\Dmin$ dependent folds.  Fixed trinomial and pentanomial formulae,
generated code, and L\'opez--Dahab reduction can be stronger under their
family assumptions~\cite{LopezDahab2000}.

Dense linear maps and generic parallel division expose a different tradeoff.
Polynomial division can be reduced to triangular Toeplitz or reciprocal
inversion~\cite{BiniPan1985}.  Barrett and Montgomery replace serial folding
with products and modulus-dependent constants.  Their machine cost cannot be
deduced from sparse shift/XOR counts.

\begin{table}[t]
\caption{Frozen direct-reduction comparison contract.  ``Primary'' denotes a
winner-eligible portable-C method; ``diagnostic'' denotes a screened baseline.}
\label{tab:families}
\centering
\small
\resizebox{\textwidth}{!}{%
\begin{tabular}{@{}lllll@{}}
\toprule
Method & Applicability & Evidence role & Reusable setup & Boundary\\
\midrule
GS & arbitrary monic & primary & doubled-shift schedule, two states & geometry/traffic\\
Serial & arbitrary monic & screened sparse baseline & tap descriptors, states & dependency chain\\
BarrettGF2X & arbitrary monic & primary & reciprocal and product buffers & gf2x thresholds\\
L\'opez--Dahab & $\deg q\le m-W$ & primary where applicable & taps and $2m$-bit buffer & restricted degree\\
Dense map & fixed, map $\le64$ MiB & diagnostic & packed $m\times m$ map & quadratic storage\\
Naive & arbitrary monic & correctness/portability & lightweight context & not tuned\\
\bottomrule
\end{tabular}
}
\end{table}

Generated fixed-modulus code is outside the general-code winner contract.
Montgomery is also excluded: REDC maps $A$ to $AR^{-1}\bmod g$, rather than
performing the frozen direct reduction $A\mapsto A\bmod g$.

\section{Sparse Feedback Operators}
\label{sec:operators}

\begin{lemma}[Low/high decomposition]
\label{lem:decomposition}
For~\eqref{eq:low-high} and~\eqref{eq:UV},
\begin{equation}
  A-gY=L+V(Y)+x^m\bigl(H+Y+U(Y)\bigr).
\end{equation}
Consequently, if $H=(I+U)Y$, then $A\bmod g=L+V(Y)$.
\end{lemma}
\begin{proof}
Expand $A-(x^m+q)Y$ and use subtraction equals addition in characteristic two.
\prooftodo
\end{proof}

\begin{lemma}[Sparse Frobenius powers]
\label{lem:frobenius}
For every $k\ge0$,
\begin{equation}
  U^{2^k}=\sum_{t\in T}N^{2^k\Delta_t},
  \label{eq:frobenius}
\end{equation}
where shifts by at least $m$ are zero.
\end{lemma}
\begin{proof}
The shifts commute, and characteristic-two Frobenius removes mixed binomial
terms.  Truncation removes shifts of distance at least $m$.
\prooftodo
\end{proof}

\begin{theorem}[Sparse feedback inverse]
\label{thm:inverse}
If $U=0$, then $(I+U)^{-1}=I$.  Otherwise, for
\begin{equation}
  r=\left\lceil\log_2\frac{m}{\Dmin}\right\rceil,
  \label{eq:r}
\end{equation}
we have $U^{2^r}=0$ and
\begin{equation}
  (I+U)^{-1}=\prod_{k=0}^{r-1}(I+U^{2^k}).
  \label{eq:inverse}
\end{equation}
\end{theorem}
\begin{proof}
Lemma~\ref{lem:frobenius} gives $U^{2^r}=0$.  Multiplication of the factors
gives $(I+U)\prod_{k<r}(I+U^{2^k})=I+U^{2^r}=I$.
\prooftodo
\end{proof}

This is a factored truncated reciprocal.  Formal-series inversion itself is
classical~\cite{Kalorkoti1993}; the proposed contribution is the sparse
fixed-length realization and its complete tap-geometry analysis.

\section{Generalized Suwako Reduction}
\label{sec:algorithm}

For $0\le k<r$, define
\begin{equation}
  T_k=\{t\in T:2^k\Delta_t<m\}.
\end{equation}
The public schedule stores shifts $2^k\Delta_t$ for $t\in T_k$ and the taps
for final low assembly.

\begin{quote}
\textbf{Algorithm 1 (generalized Suwako reduction).}
Split $A=L+x^mH$ and set $Y\leftarrow H$.  For $k=0,\ldots,r-1$, read one
immutable stage input $Y_{\rm old}$ and compute synchronously
\begin{equation*}
  Y_{\rm new}=Y_{\rm old}
  +\sum_{t\in T_k}N^{2^k\Delta_t}Y_{\rm old}.
\end{equation*}
Set $Y\leftarrow Y_{\rm new}$ and return $L+V(Y)$.
\end{quote}

An in-place tap loop that exposes updates within a stage computes a different
recurrence.

\begin{theorem}[Correctness]
\label{thm:correctness}
For every monic $g\in\poly$ of degree $m$ and every $A$ with $\deg A<2m$,
Algorithm~1 returns $A\bmod g$.
\end{theorem}
\begin{proof}
The stages apply Theorem~\ref{thm:inverse}, so
$Y=(I+U)^{-1}H$.  Lemma~\ref{lem:decomposition} then identifies $L+V(Y)$ as
the remainder.
\prooftodo
\end{proof}

\begin{example}
For $g=x^m+x^{m-d}+1$, feedback shifts are $d,2d,4d,\ldots$ while below $m$.
\end{example}

\drafttodo{Add one complete multi-tap example and cover $q=0$,
constant-only, constant-free, dense, reducible, and non-word-aligned cases.}

\section{Work--Depth--Space Geometry}
\label{sec:complexity}

Let $h_k=|T_k|$.  Algorithm~1 has exactly
$r=\lceil\log_2(m/\Dmin)\rceil$ feedback stages when $U\ne0$.  This is
feedback depth, not gate depth.

\begin{theorem}[Geometry-sensitive scheduled work]
\label{thm:work}
The feedback factors have scheduled coefficient work
\begin{equation}
  \Wfb=\sum_{t\in T}
  \sum_{k\ge0}[m-2^k\Delta_t]_+,
  \qquad [z]_+=\max\{z,0\}.
  \label{eq:work}
\end{equation}
If $s=|T|\ge1$, then
\begin{equation}
  \Wfb<m\sum_{t\in T}\left\lceil\log_2\frac{m}{\Delta_t}\right\rceil
  =O\!\left(ms\left(1+\log\frac{m}{s}\right)\right).
\end{equation}
\end{theorem}
\begin{proof}
A shift by $2^k\Delta_t$ affects $[m-2^k\Delta_t]_+$ positions.  Sum over
active taps and stages.  For the last bound, order the distinct distances and
use that the $i$th smallest positive distance is at least $i$.
\prooftodo
\end{proof}

\begin{table}[t]
\caption{Status of the mathematical claims and their completed computational
checks.  Passing tests is falsification evidence, not proof.}
\label{tab:theorem-status}
\centering
\small
\begin{tabular}{@{}p{0.27\textwidth}p{0.18\textwidth}p{0.45\textwidth}@{}}
\toprule
Statement & Status & Completed evidence\\
\midrule
Feedback decomposition and correctness & proof audit open & 20,000 random
GF(2) trials and 299,592 exhaustive binary modulus/input pairs\\
Frobenius factors and inverse & proof audit open & independent agreement over
GF(2), $\mathbb F_4$, dual numbers, and an $\mathbb F_3$ sign/radix check\\
Exact feedback-stage count & candidate theorem & native schedule checks at all
72 tested $(m,\Dmin)$ coordinates\\
Scheduled work $\Wfb$ and coarse bound & candidate theorem & all 262,125
nonempty supports for $2\le m\le18$ pass; fixed-$m$ runtime correlations are
$0.987$--$0.993$\\
Random-support laws & conjectural target & exact geometry for 10,623 frozen
supports, including median, p90, and p99\\
\bottomrule
\end{tabular}
\end{table}

Exact work depends on every tap distance, not only on $(h,\Dmin)$.  A portable
synchronous implementation can use two $m$-bit buffers and a public schedule;
generated code trades traversal for code size.

\drafttodo{Prove exact space, schedule-size, and setup-work bounds, and give a
model-consistent cost account for every ordinary reducer retained in the
evaluation.  For GS, serial folding, and L\'opez--Dahab, derive scheduled fold
and word-contribution counts; in particular, compare
$C_{\rm LD}=\Theta(\lceil m/W\rceil(h-1))$ under $\deg q<m-W$ with
$C_{\rm GS}^{\rm sched}=\Theta(\lceil m/W\rceil+
\sum_{t,k:\,2^k\Delta_t<m}\lceil(m-2^k\Delta_t)/W\rceil+
\sum_t\lceil(m-t)/W\rceil)$.  For Barrett, record the number and operand
sizes of polynomial multiplications; for the dense map, report setup, storage,
and row-parity work; analyze Naive only if its timing is retained.  Keep these
models distinct, separate aligned, cross-word, and cross-vector shifts where
applicable, and do not identify source-level counts with machine
instructions.  Either formalize the bounded-fan-in lower bound or omit it.}

\section{Prior Art and Novelty Boundary}
\label{sec:related}

Reciprocal and triangular Toeplitz inversion for polynomial division are
classical~\cite{BiniPan1985}, as are formal-power-series inversion
\cite{Kalorkoti1993} and parallel recurrence methods~\cite{KoggeStone1973}.
The inverse identity is not itself a novelty claim.  Sparse binary-field
reduction includes strong specialized methods
\cite{LopezDahab2000,NiehuesEtAl2018}; Barrett and Montgomery provide
multiplication-based alternatives, including restricted precomputation-free
families~\cite{KnezevicEtAl2008}.

Repeated Frobenius powering in factorization gives the closest application
context.  Von zur Gathen and Gerhard combine theory, implementation, storage,
and complete factorization over $\Ftwo$~\cite{GathenGerhard2002}; irreducible
polynomial construction gives another context~\cite{Shoup1990}.

Our narrow boundary is the combination of arbitrary monic-modulus
correctness, sparse support-sensitive factors, logarithmic feedback depth, no
materialized dense reciprocal or matrix, and exact geometry from all tap
distances.  We do not claim the first reciprocal, Frobenius identity, parallel
recurrence solver, or balanced XOR realization.

\drafttodo{Insert source-checked TePLAT, LFSR, CRC, prefix, and dense
linear-circuit comparisons.  Verify all metadata against primary sources.}

\section{Implementations and Evaluation}
\label{sec:evaluation}

Existing falsification suites compare with independent long division over
$\Ftwo$ and separately test $\mathbb F_4$, the nonreduced algebra
$\Ftwo[\varepsilon]/(\varepsilon^2)$, exhaustive support geometry, and an
$\mathbb F_3$ sign check.  These tests can expose errors, but do not prove
Theorem~\ref{thm:correctness}.

The native API exposes generalized Suwako, serial folding, naive long
division, and gf2x-backed Barrett under one contract.  Generalized Suwako and
serial folding share the word-shift primitive.  The ordinary-loop
L\'opez--Dahab reducer is included where its degree assumption holds, and a
dense map is screened diagnostically.  Montgomery REDC is excluded because it
computes a different Montgomery-domain operation.

\drafttodo{Record the experiment commit, commands, seeds, case counts, tap
manifests, compiler, gf2x archive, platform, affinity/frequency policy, input
distributions, timed boundaries, and setup policy.}

\placeholderbox{Predicted stage count and active-tap work versus instrumented
counts and cycles, including deviations from traffic, barriers, and alignment.}

For fixed $m$, the planned phase diagram places Hamming weight $h$ horizontally
and $m/\Dmin$ on a base-two logarithmic vertical axis.  Each support is colored
by the matched winner among the applicable GS, serial, L\'opez--Dahab, and
Barrett implementations.

\placeholderbox{Fixed-$m$ winner panels with median, p90, p99, uncertainty
marks, complete tap manifests, and friendly and hostile supports.}

\placeholderbox{Work--feedback-depth--setup Pareto map, fixed-gap and
fixed-weight sweeps, and setup amortization over reductions per modulus.}

\begin{table}[t]
\caption{Modulus provenance and the resulting claim boundary.  The timing
corpora deliberately do not infer reducibility or irreducibility from their
construction.}
\label{tab:modulus-scope}
\centering
\small
\begin{tabular}{@{}p{0.28\textwidth}p{0.20\textwidth}p{0.42\textwidth}@{}}
\toprule
Corpus & Irreducibility status & Completed evidence\\
\midrule
Controlled fixed-weight grid & unclassified & 1,096 paper-grade
reduction-only timing cells; no field-performance claim\\
Frozen random fixed-weight supports & unclassified & exact geometry for
10,623 supports; no random-support timing claim\\
Arbitrary monic correctness cases & possibly reducible & differential and
theorem-falsification checks; correctness needs only monicity\\
Named real irreducible moduli & certificate required & not yet timed;
real-modulus and field-level performance are deferred\\
\bottomrule
\end{tabular}
\end{table}

\begin{table}[t]
\caption{Portable-C reduction summary on the controlled constant-free
$\Dmin=1$ slice.  GS time and the first ratio use $h=9$.  The crossover is
the first stable Barrett winner after which every later stable sampled weight
is also Barrett; uncertain cells neither create nor break it.}
\label{tab:portable-results}
\centering
\small
\begin{tabular}{@{}rrrrr@{}}
\toprule
$m$ & GS (ns) & Barrett/GS & crossover $h$ & ratio at crossover\\
\midrule
128    & 64.1    & 2.164  & 33  & 0.881\\
512    & 272.9   & 3.807  & 97  & 0.718\\
2048   & 1079.2  & 4.247  & 97  & 0.664\\
8192   & 4390.2  & 8.654  & 129 & 0.819\\
32768  & 17008.4 & 16.768 & 257 & 0.881\\
131072 & 68167.1 & 37.712 & 641 & 0.886\\
\bottomrule
\end{tabular}
\end{table}

These medians use \texttt{portable-scalar-c:v1},
\texttt{reduction-steady-state:v1}, and the \texttt{gf2x:v1} multiplication
backend on the recorded primary machine.  They establish a platform-specific
phase boundary, not a cross-platform ordering.

\begin{table}[t]
\caption{Median setup time and requested plan-owned bytes at
$(h,\Dmin)=(9,1)$.  Each entry is ``nanoseconds / bytes.''}
\label{tab:setup-storage}
\centering
\small
\begin{tabular}{@{}rrrrr@{}}
\toprule
$m$ & trials & GS & Serial & Barrett\\
\midrule
128    & 127 & 596.5 / 744    & 119.5 / 416   & 782.0 / 184\\
512    & 127 & 719.5 / 888    & 123.5 / 512   & 4770.5 / 376\\
2048   & 127 & 887.5 / 1176   & 137.0 / 896   & 50799.0 / 1144\\
8192   & 127 & 1152.0 / 2040  & 149.0 / 2432  & 558369.5 / 4216\\
32768  & 127 & 1312.0 / 5208  & 235.0 / 8576  & 7760258.5 / 16504\\
131072 & 31  & 2832.5 / 17592 & 856.5 / 33152 & 119030566.0 / 65656\\
\bottomrule
\end{tabular}
\end{table}

Setup follows \texttt{modulus-plan:v1}.  Storage follows
\texttt{requested-owned-bytes:v1} and excludes allocator overhead, shared
modulus storage, benchmark buffers, and transient library workspace.

\section{Computational-Algebra Consequences and Limitations}
\label{sec:consequences}

Distinct-degree factorization, Rabin irreducibility testing, and sparse
polynomial search repeatedly compute Frobenius powers.  A valid experiment
must separate square formation, reduction, multiplication, and GCD cost before
recombining them in a complete workload.  A high-monomial microbenchmark is
not a square-shaped input distribution.

\placeholderbox{Square formation, reduction, and complete modular-square
breakdown for random and application-generated Frobenius-chain states.}

\drafttodo{Implement the modular-square API and compare with the strongest
general and applicable family-specific squarers.}

The primary end-to-end candidate is a complete Rabin test, distinct-degree
factorization kernel, primitive-polynomial certification, or sparse-polynomial
search.  It must report early exits, candidate distribution, irreducibility
provenance, and time fractions.  A second question asks whether, under a frozen
objective, generalized Suwako changes the platform-optimal irreducible modulus.

\begin{table}[t]
\caption{End-to-end computational-algebra results.}
\centering
\begin{tabular}{@{}llllll@{}}
\toprule
Workload & degree & modulus family & reducer & total time & reduction share\\
\midrule
\multicolumn{6}{c}{\textbf{TODO: reproducible end-to-end data}}\\
\bottomrule
\end{tabular}
\end{table}

\drafttodo{Populate only after independent state-by-state verification.  If a
reduction gain vanishes end to end, report it.  Preserve modulus candidate
sets and irreducibility certificates and state whether the winner changes.}

Dense supports, small degrees, expensive cross-word shifts, and strong fixed
reducers can make generalized Suwako unattractive.  Multiplication-based
methods may dominate once fast polynomial multiplication is effective.  The
reduction theorem does not establish a lower bound for every fixed modulus, a
constant-time implementation, or an asymptotic improvement to complete field
multiplication.

\drafttodo{Replace generic limitations with measured boundaries and threats
to validity for compiler, machine, sample size, distribution, and setup.}

\section{Conclusion}
\label{sec:conclusion}

This draft recasts binary polynomial reduction as inversion of a nilpotent
feedback operator.  Characteristic-two Frobenius powers give a factored
schedule with logarithmic feedback depth and work governed by the full sparse
tap geometry, without first constructing a dense reciprocal or matrix.

\drafttodo{Rewrite after the proof audit and experiments.  State actual
winning and losing regions and the measured workload consequence.}

\appendix

\section{Coefficient-Algebra Extension}

For a commutative characteristic-two algebra and
$U=\sum_t a_tN^{\Delta_t}$,
\begin{equation}
  U^{2^k}=\sum_t a_t^{2^k}N^{2^k\Delta_t}.
\end{equation}
Formal scheduled support and actual nonzero support can differ over a
nonreduced algebra such as $\Ftwo[\varepsilon]/(\varepsilon^2)$.

\drafttodo{Insert the fully quantified proof.  Keep the odd-characteristic
radix identity and the remainder sign $L-V(Y)$ separate.}

\section{Validation and Additional Results}

\drafttodo{Document seeds, exhaustive ranges, independent oracles, commands,
complete tables, secondary platforms, setup/storage, generated code size,
anomalies, and all losing regions.}

\begin{thebibliography}{99}
\bibitem{BiniPan1985}
D.~Bini and V.~Pan, \emph{Fast parallel polynomial division via reduction to
triangular Toeplitz matrix inversion and to polynomial inversion modulo a
power}, Inform. Process. Lett. \textbf{21} (1985), no.~2.

\bibitem{GathenGerhard2002}
J.~von zur Gathen and J.~Gerhard, \emph{Polynomial factorization over
$\mathbb F_2$}, Math. Comp. \textbf{71} (2002), no.~240, 1677--1698,
\url{https://doi.org/10.1090/S0025-5718-02-01421-7}.

\bibitem{Kalorkoti1993}
K.~S.~K. Kalorkoti, \emph{Inverting polynomials and formal power series},
SIAM J. Comput. \textbf{22} (1993), no.~3,
\url{https://doi.org/10.1137/0222037}.

\bibitem{KnezevicEtAl2008}
M.~Knezevic, K.~Sakiyama, J.~Fan, and I.~Verbauwhede, \emph{Modular reduction
in $\operatorname{GF}(2^n)$ without pre-computational phase}, WAIFI 2008,
LNCS 5130, Springer, 2008, pp.~77--87.

\bibitem{KoggeStone1973}
P.~M. Kogge and H.~S. Stone, \emph{A parallel algorithm for the efficient
solution of a general class of recurrence equations}, IEEE Trans. Comput.
\textbf{C-22} (1973), no.~8.

\bibitem{LopezDahab2000}
J.~L\'opez and R.~Dahab, \emph{High-speed software multiplication in
$\mathbb F_{2^m}$}, Progress in Cryptology---INDOCRYPT 2000, LNCS 1977,
Springer, 2000.

\bibitem{NiehuesEtAl2018}
L.~Boppre Niehues, R.~Custodio, and D.~Panario, \emph{Fast modular reduction
and squaring in $\operatorname{GF}(2^m)$}, Inform. Process. Lett.
\textbf{132} (2018), \url{https://doi.org/10.1016/j.ipl.2017.12.002}.

\bibitem{Shoup1990}
V.~Shoup, \emph{New algorithms for finding irreducible polynomials over finite
fields}, Math. Comp. \textbf{54} (1990), no.~189, 435--447,
\url{https://doi.org/10.1090/S0025-5718-1990-0993933-0}.
\end{thebibliography}

\end{document}
